\slajdid{09_3-s031}
\begin{frame}[label=09_3-s031]
\gusttekst\setlength{\abovedisplayskip}{3pt}\setlength{\belowdisplayskip}{3pt}\setlength{\abovedisplayshortskip}{2pt}\setlength{\belowdisplayshortskip}{2pt}Na identičan način možemo realizovati i složenija logička kola sa više ulaza. Na primer funkciju
\[F=\overline{D+\overline A(B+C)}\]
PDN mreža kao što smo videli pravi funkciju kod paralelno povezanih tranzistora $\overline{A+B}$ a kod redno $\overline{AB}$ odnosno najbolje bi bilo da za relaizaciju funkcije posmatramo oblik $\overline F$ odnosno
\[\overline F=D+\overline A(B+C)\]
i u tom slučaju svako ILI daje paralelnu vezu a svako I rednu vezu, dok se na ulazima tranzistora nalaze promenljive sa pravim vrednostima iz ovakvog načina prikazivanja
\par\medskip\centering\begin{tikzpicture}[x=.7cm,y=.58cm,font=\sitneoznake,>=Latex]\begin{scope}[shift={(0,1)}]\draw(0,-.35)--(0,.35);\draw(-.13,-.4)--(-.13,.4);\draw(0,.3)--(.4,.3)--(.4,.65);\draw(0,-.3)--(.4,-.3)--(.4,-.65);\draw(-1,0)node[left]{D}--(-.13,0);\end{scope}\begin{scope}[shift={(3.2,2.8)}]\draw(0,-.35)--(0,.35);\draw(-.13,-.4)--(-.13,.4);\draw(0,.3)--(.4,.3)--(.4,.65);\draw(0,-.3)--(.4,-.3)--(.4,-.65);\draw(-1,0)node[left]{$\overline A$}--(-.13,0);\end{scope}\begin{scope}[shift={(2,1)}]\draw(0,-.35)--(0,.35);\draw(-.13,-.4)--(-.13,.4);\draw(0,.3)--(.4,.3)--(.4,.65);\draw(0,-.3)--(.4,-.3)--(.4,-.65);\draw(-1,0)node[left]{B}--(-.13,0);\end{scope}\begin{scope}[shift={(4.6,1)}]\draw(0,-.35)--(0,.35);\draw(-.13,-.4)--(-.13,.4);\draw(0,.3)--(.4,.3)--(.4,.65);\draw(0,-.3)--(.4,-.3)--(.4,-.65);\draw(-1,0)node[left]{C}--(-.13,0);\end{scope}\draw(.4,1.65)--(.4,3.6)--(3.6,3.6)--(3.6,3.45);\draw(3.6,2.15)--(3.6,1.85);\draw(2.4,1.65)--(2.4,1.85)--(5,1.85)--(5,1.65);\draw(.4,.35)--(.4,.1)--(5,.1)--(5,.35);\draw(2.4,.35)--(2.4,.1);\draw(3.6,.1)--(3.6,-.4)node[ground]{};\node[below]at(2.6,-1.5){PDN mreža};\end{tikzpicture}

\end{frame}
